\chapter{Numerical method and validation}\label{ch:numerics}
The reproducible Python entry point is \texttt{python -m lambwaves.run}. The
package separates material data, determinant evaluation, root finding, modal
response, and validation. Determinants \cref{eq:matrix-S} and \cref{eq:matrix-A} are used
without tangent division. Candidate roots are bracketed, refined with Brent's
method, and checked by a scaled residual. Bulk thresholds are evaluated through
their limiting determinant rather than accepted as roots.

Continuous branch tracking is a two-stage assignment problem. Predicted frequency
from the previous two $k$ points supplies a continuation cost; normalized
eigenfield overlap
\begin{equation}C_{mn}=\left|\int_{-h}^{h}\rho\vect U_m^*(z;k_j)
\cdot\vect U_n(z;k_{j+1})\dd z\right|\label{eq:tracking-overlap}
\end{equation}
breaks close-frequency ambiguities. The current code implements frequency
continuation and exposes eigenfield-overlap tracking as pending rather than
mislabeling sorted roots. Verified outputs therefore emphasize fundamental modes
and the first isolated S-branch minimum.

Validation includes: characteristic residual; root-scan and $k$-grid refinement;
the two limits in \cref{eq:low-limits}; analytic cutoff frequencies at $k=0$;
zero initial conditions in \cref{eq:switched-solution}; and time-step refinement.
Energy/group-velocity consistency and physical modal overlaps are test targets
whose absence blocks level-B and level-C claims. The representative aluminum
constants are computational benchmark inputs, not a characterization of the
author's specimen.
